% Version 0.24. Complete finite supplier for the literal PsiBound interface.
% Source: commit 7cbf7a0c32fca1badf05931879a37c3f8a2b03b4.
\section{The finite arithmetic supplier for the odd-minus bound}
\label{sec:odd-minus-supplier}

\status{This section proves the literal \texttt{PsiBound} interface by a
finite Jacobi-sum and Kolyvagin argument.  It includes the construction of
the auxiliary classes, the selection of auxiliary primes, the finite
length descent, the Bernoulli normalization, the residual exceptional
case, and descent to the original character eigenspace.  The global
reciprocity input is the finite duality statement reconstructed in
LD2--LD4.  The remaining classical arithmetic inputs are stated explicitly
below.  No Iwasawa main conjecture or unproved Euler-system existence
statement is used in this supplier.}

\subsection{Conventions and classical inputs}

Fix an odd prime $p$, a finite extension of $\mathbb Q_p$ with ring of
integers $O$, fraction field $K$, uniformizer $\varpi$, and residue field
$k$.  The specified embedding $O\subset\overline{\mathbb Q}_p$ is kept
throughout.  Write $A=K/O$ and normalize $v=v_O$ by $v(\varpi)=1$,
with $v(0)=\infty$.  Thus $A[a]\simeq O/(a)$ and
$\operatorname{length}_O A[a]=v(a)$ for $a\ne0$.
For the source's finite DVR subalgebra $O$, this is literally its
coefficient module \texttt{QO}: the localization
$\mathbb Q_p\otimes_{\mathbb Z_p}O$ is the fraction field $K$,
and the specified quotient is $K/O$.  Indeed $O$ is a finite
torsion-free $\mathbb Z_p$-algebra and a domain, so this localization
is a finite-dimensional domain over $\mathbb Q_p$, hence a field.
Lengths of arbitrary modules initially belong to
$\mathbb N\cup\{\infty\}$.  Let $G=\operatorname{Gal}(\Omega/\mathbb Q)$
for the specified algebraic closure, and let $\chi_p:G\to O^\times$ be
the cyclotomic character.  Decomposition and inertia groups at a chosen
place above a rational prime $q$ are denoted by $D_q$ and $I_q$; whenever
several choices matter, the assertion is required at every such place.

For a continuous character $\alpha:G\to O^\times$, put
\[
 \mathscr S(\alpha)=\ker\!\left(
 H^1(G,A(\alpha))\longrightarrow
 H^1(D_p,A(\alpha))\times\prod_q H^1(I_q,A(\alpha))\right).
 \tag{OS.1}
\]
Here $A(\alpha)$ has action $g x=\alpha(g)x$.  The product notation
expresses simultaneous local triviality, and does not assert that only
finitely many local groups are nonzero.  This is the source's
$\operatorname{ShaU}$ over $\mathbb Q$.

We use the following classical theorems in their ordinary finite forms.
They are foundational inputs, rather than newly reconstructed claims.
Cyclotomic theory identifies $\operatorname{Gal}(\mathbb Q(\zeta_n)/\mathbb Q)$
with $(\mathbb Z/n)^\times$, identifies the integers with
$\mathbb Z[\zeta_n]$, gives ramification only at primes dividing $n$,
and identifies Frobenius at $q\nmid n$ with $\zeta_n\mapsto\zeta_n^q$.
At a prime $q$, roots of unity of order prime to $q$ reduce injectively;
adjoining an $M$th root of a local unit is unramified when $q\nmid M$.
The latter follows also directly from Hensel's lemma over the maximal
unramified extension, whose residue field is algebraically closed.
Chebotarev supplies Frobenius elements in every nonempty open conjugacy
class condition, outside a prescribed finite set.  Together with the
Hilbert class field and its Artin isomorphism, it supplies a degree-one
prime in every ideal class of a number field, outside a prescribed finite
set.  Indeed degree $\ge2$ primes have Dirichlet density zero: their
norm-reciprocal series converges at $s=1$, being bounded by a fixed
multiple of $\sum_r r^{-2}$ over rational primes.  The desired Artin
class has positive Dirichlet density.  Finally, Kummer theory gives
$H^1(\mathbb Q,\mu_p)=\mathbb Q^\times/\mathbb Q^{\times p}$, with its
usual compatible local identifications; this is the Kummer exact
sequence and Hilbert's theorem 90.  The exact global/local finite
duality input is isolated in OS5 and supplied by LD2--LD4.

The source map at the end of this component records pinned declarations
for each use.  In particular, these named classical inputs do not stand
for the stronger assertion that the auxiliary classes constructed below
already exist or satisfy the required recurrence.

\subsection{Jacobi elements and their two distribution laws}

Write $K_n=\mathbb Q(\zeta_n)\subset\Omega$ and
$\sigma_t(\zeta_n)=\zeta_n^t$.  For $1\le b<n$ define
\[
 \theta_{n,b}(t)=
 \left\lfloor\frac{b\langle t\rangle_n}{n}\right\rfloor-
 \left\lfloor\frac{b\langle-t\rangle_n}{n}\right\rfloor,
 \qquad t\in(\mathbb Z/n)^\times.
 \tag{OS.2}
\]
Thus $\theta_{n,b}(-t)=-\theta_{n,b}(t)$.
For an ideal $\mathfrak a\ne0$ of $K_n$, call $X$ an associated
Jacobi element if
\[
 X\overline X=1,\qquad
 (X)=\prod_{t\in(\mathbb Z/n)^\times}
       \sigma_t^{-1}\mathfrak a^{\,\theta_{n,b}(t)}.
 \tag{OS.3}
\]
The inverse in this formula fixes the convention used in all subsequent
Frobenius and character formulas.

% BEGIN OS1
\begin{proposition}[OS1: existence, norm, and congruence of Jacobi elements]
For every nonzero ideal $\mathfrak a$ an element satisfying (OS.3)
exists, and any two differ by a root of unity whose order divides $2n$.
Let $q\nmid n$ be prime, $b<q$, and let $X'$ and $X$ satisfy (OS.3)
at levels $nq$ and $n$, for $\mathfrak a'$ and
$N_{K_{nq}/K_n}\mathfrak a'$, respectively.  Then
\[
 N_{K_{nq}/K_n}X'
   =\zeta\,X/\sigma_q^{-1}X,\qquad \zeta^{2n}=1.
 \tag{OS.4}
\]
If $\mathfrak a'$ is prime to $q$, there is a single root of unity
$\xi$, with $\xi^{2nq}=1$, such that, at every prime above $q$,
\[
                  X'\equiv\xi\,\sigma_q^{-1}X\pmod q.
 \tag{OS.5}
\]
Congruences here mean equality in the local residue field; all the
elements being compared are local units.
\end{proposition}
\begin{proof}
First suppose $\mathfrak a=P$ has prime norm $r\nmid n$.
Then $r\equiv1\pmod n$, and reduction at $P$ identifies
$\mu_n$ with the subgroup of order $n$ in $\mathbb F_r^\times$.
Let $d=(r-1)/n$ and let $\eta:\mathbb F_r^\times\to\mu_n$ be
the character reducing to $x\mapsto x^d$, extended by zero at zero.
Set
\[
 J=\prod_{j=1}^{b-1}J(\eta,\eta^j),\quad
 J(\eta,\eta^j)=\sum_{x\in\mathbb F_r}\eta(x)\eta^j(1-x),
 \qquad X=\overline J/J.
 \tag{OS.6}
\]
The standard finite-sum change of variables gives
$J(\eta,\eta^j)J(\eta^{-1},\eta^{-j})=r$ when
$\eta,\eta^j,\eta^{j+1}$ are nontrivial.  One way to check the
identity is to introduce a nontrivial additive character, expand the
two Gauss sums, and substitute $x+y=u$; this gives
$g(\eta)g(\eta^j)=J(\eta,\eta^j)g(\eta^{j+1})$ and
$g(\lambda)g(\lambda^{-1})=\lambda(-1)r$.
Consequently each Jacobi sum is nonzero and its integral ideal divides
$(r)$.

Put $u=\langle t\rangle_n$ and $v=\langle jt\rangle_n$.
Reduction at $\sigma_t^{-1}P$ gives
\[
 J(\eta,\eta^j)\equiv
       \sum_{x\in\mathbb F_r}x^{du}(1-x)^{dv}.
\]
For $0<i,j<r-1$, the sum $\sum_x x^i(1-x)^j$ vanishes exactly
when $i+j<r-1$.  To see this, expand $(1-x)^j$; the only possible
nonzero power sum has exponent $r-1$, and its coefficient is
$-(-1)^{r-1-i}\binom{j}{r-1-i}$, which is nonzero modulo $r$ when
$0\le r-1-i\le j<r$.  Thus the exponent of
$\sigma_t^{-1}P$ in the Jacobi sum is
$1-\operatorname{carry}(u,v)$, where the carry is $0$ or $1$
according as $u+v<n$ or $u+v\ge n$.
The prime $r$ splits without ramification, so the exponents really are
zero or one.  Telescoping carries gives
\[
 \operatorname{ord}_{\sigma_t^{-1}P}J
    =b-1-\left\lfloor\frac{b\langle t\rangle_n}{n}\right\rfloor.
\]
Subtracting this formula from its conjugate proves (OS.3) for
$\overline J/J$.

For a general $\mathfrak a$, choose a degree-one prime $P$, prime to
$n$, whose ideal class is $[\mathfrak a]^{-1}$, and write
$\mathfrak aP=(y)$.  The explicit element
\[
                 Y=\prod_t\sigma_t^{-1}(y)^{\theta_{n,b}(t)}
 \tag{OS.7}
\]
satisfies (OS.3) for $(y)$: the ideal identity is immediate and the
conjugate factors cancel in pairs.  Dividing $Y$ by the element (OS.6)
for $P$ proves existence for $\mathfrak a$.
The quotient of two choices is a unit all of whose complex conjugates
have absolute value one.  It is a root of unity: the monic polynomials
whose roots are the conjugates of its powers have integral, uniformly
bounded coefficients, hence only finitely many possibilities; some two
powers agree.  If a root of order $h$ lies in $K_n$, then
$K_{\operatorname{lcm}(n,h)}=K_n$, so their Euler totients are equal.
The formula for the totient permits no increase in a prime-power
factor, except adjoining a factor $2$ when $n$ is odd.  Thus
$h\mid2n$.  This proves
the stated uniqueness.

The norm identity rests on the following integer identity:
\[
 \sum_{\substack{t'\in(\mathbb Z/nq)^\times\\t'\equiv t\pmod n}}
 \theta_{nq,b}(t')
       =\theta_{n,b}(t)-\theta_{n,b}(q^{-1}t).
 \tag{OS.8}
\]
To verify it, first sum over the $q$ representatives $u+in$ of $t$,
$0\le i<q$.  Since $0\le b<q$, the usual floor-sum identity gives
the floor at level $n$, plus the constant
$\sum_{i=0}^{q-1}\lfloor ib/q\rfloor$.
Remove the unique representative divisible by $q$; its contribution
is the floor at $q^{-1}t$.  Subtract the identical computation for
$-t$; the constants cancel.  Apply (OS.8) to ideals and use norm
compatibility with conjugation.  The two sides of (OS.4), with the
root of unity omitted, then have the same ideal and conjugate product
one.  The uniqueness just proved supplies $\zeta$.

For (OS.5), first take $\mathfrak a'$ of degree one and norm $r$.
The characters $\eta'$ at level $nq$ and $\eta$ at level $n$ in
(OS.6) satisfy $(\eta')^q=\eta$ after inclusion in the larger field:
both sides have the same reduction at the prime above $r$, and
reduction is injective on $\mu_{nq}$.  Also
$(\sigma_q^{-1}\eta)^q=\eta$.  In characteristic $q$, the Frobenius
map is injective, so $\eta'$ and $\sigma_q^{-1}\eta$ have the same
reduction at every prime above $q$.  The Jacobi sums, their products,
and their conjugates therefore have the same reductions.  Their
denominators are units there, since they divide powers of $r\ne q$.
This proves (OS.5) for the canonical choices, and uniqueness inserts
one root of unity for any other choices.

For a principal ideal $(y)$ prime to $q$, use (OS.7).
Reduction of $\sigma_{t'}^{-1}y$ depends only on $t'\bmod n$,
because $\zeta_q$ reduces to one.  Call this value $f(t)$.
Frobenius gives $f(t)=f(qt)^q$.  Combining this with (OS.8), the
reduction of the upper-level product is
\[
 \prod_t f(t)^{\theta(t)-\theta(q^{-1}t)}
   =\prod_t f(qt)^{(q-1)\theta(t)}.
\]
Every one of the $q-1$ relative conjugates of $y$ has the same
reduction.  The last expression is consequently the reduction of
the lower-level product attached to $N(y)$, acted on by
$\sigma_q^{-1}$.  This proves (OS.5) for principal ideals.
Finally choose a degree-one prime $P'$ in $[\mathfrak a']^{-1}$,
outside $nq$, and apply the principal calculation to
$P'\mathfrak a'=(y)$.  Divide by the already proved calculation
for $P'$.  All factors are units at $q$, and the root of unity from
the degree-one calculation is global, so a single $\xi$ works at
every prime above $q$.  This proves (OS.5) in full.
\end{proof}
% END OS1

\subsection{Removing the roots of unity}

For the ensuing finite construction let $M=p^m$, $m\ge1$, and
$F=N_0M$ with $(N_0,p)=1$.  Suppose
\[
 \alpha|_{U}\equiv1\pmod M,\quad
 \alpha(j)=-1,\quad
 \alpha(g_0)-\chi_p(g_0)\in O^\times,
 \qquad U=\operatorname{Gal}(\Omega/K_F),
 \tag{OS.9}
\]
where $j$ is a complex conjugation.  Put
$\beta=\chi_p\alpha^{-1}$.  Thus $\beta(j)=1$,
$\beta|_U\equiv1\pmod M$, and some $\beta(g)-1$ is a unit.
Choose an integer $a$ such that $g_0$ acts on $\mu_M$ by the
$a$th power, and put
\[
                   P(y)=g_0(y)/y^a.
 \tag{OS.10}
\]
Then $\alpha(g_0)-a$ is a unit, since $M$ is in the maximal ideal.

Choose a prime $r\equiv1\pmod F$ and one place $A_0$ above $r$.
Call a prime $q$ admissible if
\[
 q\equiv1\pmod F,\qquad q\ne r,\qquad
              p\nmid\operatorname{ord}_{\mathbb F_q^\times}(r).
 \tag{OS.11}
\]
An admissible list has distinct admissible entries, and $d$ denotes
their product.  All cyclotomic fields below are the actual subfields
of the chosen $\Omega$.

% BEGIN OS2
\begin{proposition}[OS2: a strict finite Euler system]
For $1\le b<F$ there are nonzero elements $x_d\in K_{Fd}$,
units outside $r$, and positive integers $s(d,q)$ prime to $p$,
such that for every admissible extension $d\mapsto dq$,
\begin{align*}
 N_{K_{Fdq}/K_{Fd}}x_{dq}
       &=(x_d/\sigma_q^{-1}x_d)^{s(d,q)},\\
 x_{dq}&\equiv(\sigma_q^{-1}x_d)^{s(d,q)}\pmod q.
\end{align*}
The second equality holds at every place above $q$.
The elements can be chosen in the form
$x_d=P(X_d)^{2N_0d}$, where $X_1$ has
\[
 \operatorname{ord}_{A_0}(\sigma_tX_1)
                 =\theta_{F,b}(t)\operatorname{ord}_{A_0}(r).
 \tag{OS.12}
\]
\end{proposition}
\begin{proof}
At each level take the prime below the fixed place $A_0$, and use OS1
to choose its Jacobi element $X_d$.  These elements are units outside
$r$.  Relative norm of the upper prime is a power of the lower
prime, with exponent $s_0(d,q)$ dividing the order of $r$ modulo $q$.
Indeed residue degrees in cyclotomic fields are orders of $r$ modulo
their levels, and their relative ratio divides that order.  The
admissibility condition makes $s_0(d,q)$ prime to $p$.
OS1 therefore gives the displayed norm and congruence for $X_d$
with exponent $s_0(d,q)$ and possible root-of-unity factors of order
dividing $2Fdq$.  At level $F$, $r$ splits completely, so the ideal
identity (OS.3) gives precisely (OS.12).

The operator $P$ commutes with cyclotomic conjugations and with their
norms.  If $\zeta^{2Fdq}=1$, then
$\zeta^{2N_0dq}\in\mu_M$, and consequently
\[
                     P(\zeta)^{2N_0dq}=1.
\]
Apply $P$ and then the power $2N_0dq$ to both lax identities.
The root-of-unity factors disappear.  On the lower side this power
is $q$ times the exponent used in $x_d$, so the new exponent is
$s(d,q)=q s_0(d,q)$, still prime to $p$.
Congruences remain valid after applying $g_0$, because the original
congruence held at every prime above $q$, and after division and
integer powers because all relevant elements are local units.
This proves both strict identities simultaneously.
\end{proof}
% END OS2

\subsection{Derivative classes with exact local recurrence}

For every admissible $q$, choose $\tau_q$ acting trivially on all
roots of unity of order prime to $q$, and generating the action on
$\mu_q$.  Such an element exists in inertia at $q$.
In multiplicative notation put
\[
 D_q y=\prod_{i=0}^{q-2}\tau_q^i(y)^i,
 \qquad D_L=\prod_{q\in L}D_q,
 \qquad y_L=D_Lx_d.
 \tag{OS.13}
\]
The last product means composition of commuting group-ring operators
on the indicated cyclotomic field.  Fix a primitive $M$th root
$\zeta\in\Omega$, representatives $\sigma_t\in G$ of
$t\in(\mathbb Z/F)^\times$, and $w\in A[M]$.

% BEGIN OS3
\begin{proposition}[OS3: construction of the auxiliary homomorphisms]
There are continuous additive homomorphisms $\nu_L:U\to A[M]$
such that
\[
 \nu_L(gug^{-1})=\beta(g)\nu_L(u),\qquad
 \nu_L(I_q)=0\quad(q\nmid Frd).
 \tag{OS.14}
\]
If $q$ is admissible and absent from $L$, then for every place above
$q$ and every Frobenius $\phi_q$ there are
$\tau\in I_q$ and $u_q\in O^\times$ with
\[
                  \nu_{qL}(\tau)=u_q\nu_L(\phi_q).
 \tag{OS.15}
\]
At the fixed place above $r$ there are $\tau_0\in I_r$ and
$u_0\in O^\times$ with
\[
 \nu_\varnothing(\tau_0)=u_0zw,\qquad
 z=\sum_{t\in(\mathbb Z/F)^\times}
          \theta_{F,b}(t)\alpha(\sigma_t)^{-1}.
 \tag{OS.16}
\]
These are the full source conditions \texttt{NuSys}, for the same
$M,F,r,\alpha,A_0,\sigma$ and $w$.
\end{proposition}
\begin{proof}
The group-ring identity
\[
                 (\tau_q-1)D_q=(q-1)-N_q
 \tag{OS.17}
\]
and the norm formula of OS2 imply, by induction on $L$, that
$g y_L/y_L$ is an $M$th power in $K_{Fd}$ for every $g\in U$.
Here are the induction details.  For a new $q$ write
$w'=D_Lx_{dq}$ and $y=D_Lx_d$.  Then
$N_qw'=(y/\phi_q^{-1}y)^s$, where $s=s(d,q)$.
By induction choose $z'$ in $K_{Fd}$ with
$(z')^M=\phi_q^{-1}y/y$.  Since $q-1=Mh$, (OS.17) gives
\[
 \frac{\tau_q(D_qw')}{D_qw'}
          =\frac{(w')^{q-1}}{N_qw'}
          =((w')^h(z')^s)^M.
\]
Repeat this calculation for each member of the enlarged list, using
the derivative for the list with that member removed; the induction
hypothesis applies to that shorter list.  For elements fixing
$K_{Fdq}$ the quotient is one.
These elements generate $\operatorname{Gal}(K_{Fdq}/K_F)$, by the
Chinese remainder theorem.  The property is closed under products
and inverses, proving the induction.

Choose $Y_L^M=y_L$.  For $g\in U$ choose $z_g\in K_{Fd}$ with
$z_g^M=g y_L/y_L$, and define $\kappa_L(g)\in\mathbb Z/M$ by
\[
 \zeta^{\kappa_L(g)}=
       \frac{g(PY_L)}{PY_L}\,\frac1{Pz_g}.
 \tag{OS.18}
\]
The $M$th power of the right side is one, because $g_0$ and $g$
commute on $K_{Fd}$.  Changing $z_g$ multiplies it by an element of
$\mu_M$, which $P$ kills; changing $Y_L$ has the same harmless
effect.  For $g,h\in U$, one can take $z_{gh}=g(z_h)z_g$.
Substitution in (OS.18), using the same commutation on $K_{Fd}$,
proves additivity.  Its kernel contains the open subgroup fixing
$K_{Fd}$ and $PY_L$, so $\kappa_L$ is continuous.
At a prime outside $Frd$, the field $K_{Fd}$ is unramified and
$y_L$ is a unit.  Its $M$th roots are unramified there, as are their
$g_0$-conjugates.  Inertia therefore gives zero in (OS.18).

Define
\[
 \nu_L(u)=\sum_t\kappa_L(\sigma_t^{-1}u\sigma_t)
                              \beta(\sigma_t)w.
 \tag{OS.19}
\]
Additivity, continuity, $M$-torsion and the stated unramifiedness are
immediate.  To check equivariance, reindex the cosets by $g$.
The extra factors are in $U$; an additive homomorphism on $U$ is
unchanged by conjugation inside $U$, and $\beta|_U$ acts trivially
on $A[M]$.  The remaining factor is $\beta(g)$, proving (OS.14).

For completeness the local recurrence in (OS.15) is an exact residue
calculation, not an additional Euler-system axiom.  Choose inertia
$\tau$ with the same action on $\mu_q$ as $\tau_q$.
Put $w'=D_Lx_{dq}$, $y=D_Lx_d$, $Y^M=D_qw'$ and $(Y')^M=y$.
For $\phi=\phi_q$ choose $z'$ with $(z')^M=\phi^{-1}y/y$.
The previous induction shows that $(w')^h(z')^s$ is an allowable
root for $\tau$, where $Mh=q-1$.
Substituting these roots in (OS.18) and cancelling factors gives
\[
 \frac{\zeta^{\kappa_{qL}(\tau)}}
      {(\zeta^{\kappa_L(\phi^{-1})})^s}
 =\frac{\tau(PY)}{PY}
   \left(\frac{PY'}{\phi^{-1}(PY')^q}\right)^s
   \left[P\!\left(\frac{w'}{(\phi^{-1}y)^s}\right)\right]^{-h}.
 \tag{OS.20}
\]
The first factor reduces to one because $\tau$ is inertia; the
second does so because $\phi$ is Frobenius; the third does so by
the congruence in OS2 and its images under $D_L$ and $P$.
All elements involved are units at $q$.
Both sides on the left belong to $\mu_M$, whose reduction is
injective since $q\ne p$.  Hence
$\kappa_{qL}(\tau)=-s\kappa_L(\phi)$ in $\mathbb Z/M$.
The same identity holds after conjugating the chosen place by each
$\sigma_t$: cyclotomic automorphisms commute, so the action of
the conjugated inertia element on $\mu_q$ is still the required
generator.  Sum as in (OS.19).  This gives (OS.15) with
$u_q=-s$, a unit.

For initialization, if $y\in K_F$ and
$\operatorname{ord}_{A_0}(\sigma_tPy)=k(t)
\operatorname{ord}_{A_0}(r)$, choose $\rho^M=r$ and tame inertia
$\tau_0$ with $\tau_0\rho=\zeta\rho$.
The quotient $\sigma_t(PY)/\rho^{k(t)}$ is a local unit whose
$M$th power is a unit, so it is fixed by inertia.  Formula (OS.18)
therefore gives
\[
       t\,\kappa(\sigma_t^{-1}\tau_0\sigma_t)=k(t)
                         \quad\hbox{in }\mathbb Z/M.
\]
Since $\chi_p(\sigma_t)\equiv t\pmod M$, its weighted sum is
$\sum_t k(t)\alpha(\sigma_t)^{-1}w$.
For $y=x_1=P(X_1)^{2N_0}$, write $d_0$ for the image of $g_0$
in $(\mathbb Z/F)^\times$.  By (OS.12),
\[
 k(t)=2N_0\bigl(\theta(td_0^2)-2a\theta(td_0)+a^2\theta(t)\bigr).
\]
Reindexing the sum, and using $\alpha|_U\equiv1\pmod M$, gives
\[
 \sum_t k(t)\alpha(\sigma_t)^{-1}w
       =2N_0(\alpha(g_0)-a)^2zw.
\]
The coefficient is a unit by (OS.9) and $p\ne2$.
This proves (OS.16) and the entire assertion.
\end{proof}
% END OS3

\subsection{Compatible auxiliary primes}

% BEGIN OS4
\begin{proposition}[OS4: one prime with all the required properties]
Assume (OS.9).  Let $c:U\to A[M]$ be a nonzero continuous
$\alpha$-equivariant homomorphism, and let
$\nu:U\to A[M]$ be a continuous $\beta$-equivariant homomorphism.
For every finite set of excluded primes, there are an admissible prime
$q$ outside that set, a place above $q$, and its Frobenius $\phi$, such
that
\[
       c(\phi)\ne0,\qquad
       O\nu(\phi)=\sum_{u\in U}O\nu(u).
 \tag{OS.21}
\]
The same prime can in addition be required to satisfy
$q\not\equiv1\pmod{Fp}$.
\end{proposition}
\begin{proof}
Choose $\rho^M=r$ and put $H=\operatorname{Gal}(\Omega/K_F(\rho))$.
The extension $K_F(\rho)/K_F$ is cyclic of degree $M$:
$X^M-r$ is Eisenstein at a prime above $r$, and $\mu_M\subset K_F$.
It is Galois over $\mathbb Q$, because every conjugate of $\rho$
is $\zeta^i\rho$.  On its relative cyclic group, conjugation by
$g\in G$ acts by the scalar $\chi_p(g)$ modulo $M$.

Let $B=\ker c$.  If $H\subset B$, then $c$ factors through this
cyclic Kummer group.  Comparing its $\alpha$-equivariance with
the preceding conjugation action gives
$(\alpha(g)-\chi_p(g))c(u)=0$ for every $g,u$.
At $g_0$ the coefficient is a unit, contrary to $c\ne0$.
Thus $H\not\subset B$.

The $O$-module $C=\sum O\nu(U)$ is finite cyclic: $\nu$ has finite
image by compactness, and every finite submodule of $A$ is one of
the modules $A[\varpi^a]$.
If $C\ne0$, let $B'=\{u:\nu(u)\in\varpi C\}$.
If $H\subset B'$, the nonzero homomorphism
$U\to C/\varpi C$ would again factor through the Kummer group.
Its two conjugation scalars would be $\beta(j)=1$ and
$\chi_p(j)=-1$.  Their difference is a unit because $p$ is odd,
a contradiction.  Thus $H\not\subset B'$.
If $C=0$, simply set $B'=B$; the second condition of (OS.21) is
then automatic.

Put $B''=\operatorname{Gal}(\Omega/K_{Fp})$, a subgroup of $U$.
We also have $H\not\subset B''$.
Otherwise $K_{Fp}\subset K_F(\rho)$.  The relative extension
$K_{Fp}/K_F$ has degree $p$ and is abelian over $\mathbb Q$.
Complex conjugation would therefore act trivially on its relative
Galois group.  But, as a quotient of the Kummer group, that action
is inversion.  A nontrivial group of odd order $p$ cannot have
both actions.  This proves the third noncontainment.

Here is the small group-theoretic compatibility argument.
A group cannot be the union of three proper subgroups $A_1,B_1,C_1$
when $x^M\in B_1$ for every $x$ and $M$ is odd.
Two proper subgroups cannot cover a group: the product of an element
in each exclusive part lies outside both.
If a three-subgroup cover existed, choose
$x\notin B_1\cup C_1$ and $y\notin A_1\cup C_1$ using this
two-subgroup fact.  Then $x\in A_1$, $y\in B_1$, and $xy\in C_1$.
The element $x^2y$ is in neither $A_1$ nor $C_1$, so belongs to
$B_1$.  Hence $x^2\in B_1$.  Combining this with $x^M\in B_1$
and oddness of $M$ gives $x\in B_1$, a contradiction.
Apply the argument inside $H$ to its intersections with $B',B,B''$.
The required power condition follows because $c(x^M)=Mc(x)=0$.
Hence there is an element
\[
                         h\in H\setminus(B\cup B'\cup B'').
\]
All three subgroups are open in $U$; thus this is a nonempty open
condition in $G$.  Pass to a finite Galois quotient small enough
to preserve it.  Chebotarev supplies a Frobenius conjugate in this
open set outside the prescribed primes.  Changing the place above
that rational prime selects the indicated conjugate itself.
This explains compatibility of all conditions in the compositum;
no independent, possibly conflicting congruence choices have been
made.

The selected Frobenius lies in $U\setminus B''$, so
$q\equiv1\pmod F$ but $q\not\equiv1\pmod{Fp}$.
Write $q-1=N_0M h_0$; then $p\nmid h_0$.
It fixes $\rho$, so its residue satisfies $\rho^{q-1}=1$,
and consequently $r^{N_0h_0}=1$ modulo $q$.
Therefore $\operatorname{ord}_{\mathbb F_q^\times}(r)$ divides
$N_0h_0$, which is prime to $p$.  Excluding $r$ makes $q$
admissible.  Avoidance of $B$ gives the first assertion of (OS.21),
and avoidance of $B'$ says that $\nu(\phi)$ is a generator of
the cyclic module $C$.  This proves every required condition.
\end{proof}
% END OS4

\subsection{The precise reciprocity consequence}

% BEGIN OS5
\begin{proposition}[OS5: finite reciprocity and the annihilator comparison]
The duality results LD2--LD4 imply the source condition
\texttt{Rec}.  In particular, let $a_0\mid a\mid a'$ with
$a'\mid M$, and let $c:G\to A$ be an $\alpha$-cocycle killed by
$a_0$.  Assume $c$ is zero on inertia away $F$, is a coboundary
on decomposition groups above $F$ with witnesses killed by $a_0$,
and is zero on the decomposition groups above the primes dividing
$n$ outside $F$.
Suppose $q\equiv1\pmod F$, $q\nmid Fn$, and
$\nu:U\to A[M]$ is $\beta$-equivariant, unramified outside $Fnq$,
with
\[
 \operatorname{Ann}_O\nu(U)=(a'),\qquad
                  \operatorname{Ann}_O\nu(\tau)=(a)
 \tag{OS.22}
\]
for some $\tau\in I_q$.
Then $a'=a d$ for some $d\in O$ and
\[
                         d\,c(\phi_q)=0.
 \tag{OS.23}
\]
\end{proposition}
\begin{proof}
We first spell out the finite coefficient reduction, so that the
reference to LD4 preserves the literal \texttt{Rec} quantifiers.
For any nonzero annihilator $a'$ choose $w_0\in A$ with
$\operatorname{Ann}(w_0)=(a')$ and identify $A[a']$ with
$B=O/(a')$ by $b\mapsto bw_0$.
Choose $p^h$ divisible by $a'$ and let
$\lambda:B\to\mathbb Z/p^h$ be any
$\mathbb Z/p^h$-linear map.
For an $\alpha$-cocycle $c$ and a
$\beta=\chi_p\alpha^{-1}$-cocycle $n$, the function
\[
             (g,h)\longmapsto
                 \lambda\bigl(n(g)\alpha(g)c(h)\bigr)
 \tag{OS.24}
\]
is their cup-product cocycle in $(\mathbb Z/p^h)(1)$.
Use the assumed continuity of the reduction of $\alpha$ modulo $a'$;
this is exactly the open-kernel hypothesis in \texttt{Rec}.
No continuity of an unreduced character is needed at this step.
At primes outside the finite ramification
set this cup class is unramified; at every specified prime where
$c$ is a coboundary, it is a coboundary.  For example if
$c(g)=(\alpha(g)-1)m$, a cobounding cochain is
$-\lambda(n(g)\alpha(g)m)$.
At the single remaining prime $q$, LD4 first makes the local cup
class zero.  The following argument converts that fact into the
elementwise assertion actually needed.  Let
$x_0:D_q/I_q\to R=\mathbb Z/p^h$ send arithmetic Frobenius to $1$.
Since $c$ is unramified and the characters are trivial on $B$,
$c(g)=x_0(g)c(\phi_q)$ in $B$.
The function $N(g)=\lambda(n(g)c(\phi_q))$ is a cocycle in $R(1)$:
on these values the action of $\beta$ agrees with that of $\chi_p$,
as both are trivial modulo $a'$ on $D_q$.
Thus the local cup (OS.24) is $N\cup x_0$.
Every unramified $R$-valued class is an $R$-multiple of $x_0$,
so $N$ pairs trivially with every such class.
LD2, applied to the trivial module $R$ and its dual $R(1)$, makes
$N$ unramified.  Inertia acts trivially on $R(1)$ at $q\ne p$,
so an unramified class has $N(\tau)=0$ as an actual cocycle value.
Multiplying by $x_0(g)$ now gives
\[
                  \lambda\bigl(n(\tau)c(g)\bigr)=0
                 \quad(\tau\in I_q,\ g\in D_q),
 \tag{OS.25}
\]
under exactly the stated trivial-character and unramified conditions.
This deduction uses the same finite modules and the same place choices
throughout.
All the maps $\lambda$ together separate points of $B$:
decompose its finite additive $p$-group into cyclic factors and
map a nonzero coordinate into $\mathbb Z/p^h$.
Thus (OS.25) holds before applying $\lambda$ as well.
Equivalently, $n(\tau)=yw_0$ implies $y c(g)=0$ in $A$.
This is exactly \texttt{Rec}; no trace pairing of potentially
nonunit discriminant has been substituted for these separating maps.

To apply it to $\nu$, choose $g_1$ with $\beta(g_1)-1$ a unit.
The quotient $G/U$ is abelian, so
\[
 n(g)=-(\beta(g_1)-1)^{-1}
                \nu(gg_1g^{-1}g_1^{-1})
 \tag{OS.26}
\]
is defined.  The commutator identity and $\beta$-equivariance give
$n(gh)=n(g)+\beta(g)n(h)$, and for $u\in U$ give $n(u)=\nu(u)$.
Its values are $O$-multiples of values of $\nu$, so $a'n=0$.
It is continuous, and it has exactly the unramifiedness needed
outside $Fnq$, since inertia there lies in $U$.
At $q$, $D_q\subset U$ because $q$ splits in $K_F$; hence
$\alpha,\beta\equiv1\pmod{a'}$.  The hypotheses on $c$
provide the other local conditions in \texttt{Rec}.
We obtain $y c(\phi_q)=0$ if $\nu(\tau)=yw_0$.

The inclusion $\nu(\tau)\in\nu(U)$ shows $a\mid a'$.
In a DVR,
\[
 v(a)=\max\{v(a')-v(y),0\},\qquad
 v(d)=v(a')-v(a)=\min\{v(y),v(a')\}.
\]
If $v(y)\le v(a')$, then $d$ is a multiple of $y$ up to a unit,
so it kills $c(\phi_q)$.  If $v(y)\ge v(a')$, then $a$ is a
unit, and $a_0\mid a$ makes $c=0$.
This proves (OS.23).
\end{proof}
% END OS5

\subsection{The finite descent actually bounds length}

For $\alpha(j)=-1$, let $\mathcal T(F,a_0)$ consist of continuous
$\alpha$-cocycles $c:G\to A$ with $c(j)=0$, killed by $a_0$,
zero on inertia away $F$, and locally coboundaries at primes
dividing $F$ with witnesses killed by $a_0$.
This is the exact normalized module \texttt{ShT}.
Let $\mathcal T_n$ be its submodule whose cocycles vanish on every
decomposition group above a prime dividing $n$.

% BEGIN OS6
\begin{proposition}[OS6: the complete finite length descent]
Under (OS.9), and with $z$ as in (OS.16), if
$a_0z\mid M$, then
\[
              \operatorname{length}_O\mathcal T(F,a_0)\le v(z).
 \tag{OS.27}
\]
No finiteness assumption on the module on the left is required.
\end{proposition}
\begin{proof}
The divisibility implies $z\ne0$; write $M=k_0z$, so $a_0\mid k_0$.
Choose $w\in A$ with annihilator $(M)$, and construct $\nu_L$
by OS3.  Let $(a_L)=\operatorname{Ann}\nu_L(U)$.
Every $a_L$ divides $M$, and $a_L\ne0$.
Initialization gives $\operatorname{Ann}\nu_\varnothing(\tau_0)=(k_0)$,
so $a_\varnothing=k_0d_0$.

We will repeatedly use the following length comparison, which is
stronger than an elementwise annihilation statement.
If $\nu$ has annihilator $(a')$ and $\nu(\tau)$ has annihilator
$(a)$, with $a_0\mid a$ and $a'=ad$, OS5 shows that evaluation
\[
             \mathcal T_n\longrightarrow A,\qquad
                       c\longmapsto c(\phi_q)
\]
has image in $A[d]$, of length at most $v(d)$.
Its kernel is $\mathcal T_{nq}$.
Indeed inertia at $q$ acts trivially on the values and $c$ vanishes
there; the unramified quotient is procyclic, so a continuous
homomorphism zero at Frobenius vanishes on $D_q$.
The same conclusion holds at every conjugate place: the cocycle
conjugation formula changes the value only by
$(\alpha(g)-1)c(h)$ with $g\in D_q\subset U$, which is zero
because $a_0\mid M$ and $\alpha|_U\equiv1\pmod M$.
Thus additivity and monotonicity of extended module length give
\[
        \operatorname{length}\mathcal T_n
       \le\operatorname{length}\mathcal T_{nq}+v(d).
 \tag{OS.28}
\]
Apply this first with $n=1$, $q=r$, and the initialization data:
\[
 \operatorname{length}\mathcal T(F,a_0)
       \le\operatorname{length}\mathcal T_r+v(d_0).
 \tag{OS.29}
\]

For an admissible list $L$ with $a_0\mid a_L$, put
$S_L=\mathcal T_{r\prod L}$ and $J=v(M)$.
If $S_L\ne0$, choose a nonzero $c\in S_L$.
Its restriction to $U$ is an $\alpha$-equivariant homomorphism:
the action is trivial modulo $M$, and the cocycle conjugation
formula gives equivariance.  This restriction is nonzero.
For if it were zero, $c$ would factor through the abelian group
$G/U$; the commutativity relation for $j$ and $g$, together with
$c(j)=0$, would give $(\alpha(j)-1)c(g)=0$, hence $c=0$.

Use OS4 with this $c$ and $\nu_L$, excluding primes dividing
$Fr\prod L$.  The resulting $q$ has
$c(\phi_q)\ne0$ and $\nu_L(\phi_q)$ generates $\nu_L(U)$.
By OS3, $\nu_{qL}(\tau)$ therefore has annihilator $(a_L)$.
Write $a_{qL}=a_Ld$ and apply (OS.28):
\[
 \operatorname{length}S_L+v(a_L)
       \le\operatorname{length}S_{qL}+v(a_{qL}).
 \tag{OS.30}
\]
Moreover $v(d)>0$, because (OS.23) kills the nonzero value
$c(\phi_q)$; a unit cannot do so.  Thus $v(a_{qL})>v(a_L)$,
while both are at most the finite integer $J$.

Induction on $J-v(a_L)$ now proves
\[
                   \operatorname{length}S_L+v(a_L)\le J.
 \tag{OS.31}
\]
If $S_L=0$ this is immediate.  Otherwise (OS.30) reduces it to
a strictly smaller induction parameter.  This proof also establishes
finiteness; it has not subtracted an infinite length or presumed
that a nonzero element exhausts a module.
At $L=\varnothing$, substitute
$a_\varnothing=k_0d_0$ and $M=k_0z$ in (OS.31):
\[
                  \operatorname{length}\mathcal T_r+v(d_0)\le v(z).
\]
Only the finite integer $v(k_0)$ has been cancelled.
Combine this with (OS.29) to obtain (OS.27).
\end{proof}
% END OS6

\subsection{The Bernoulli value with the correct denominator and sign}

For a primitive Dirichlet character $\eta$ of conductor $f$, set
\[
 B_{k,\eta}=f^{k-1}\sum_{t=1}^{f}\eta(t)B_k(t/f),
 \qquad \operatorname{den}(N,k)=N\prod_{\substack{\ell\le k+1\\
                                    \ell\ {
m prime}}}\ell.
 \tag{OS.32}
\]
The generalized Bernoulli number in the source's
$\operatorname{Bprim}$ is always taken at the primitive conductor.
Thus for an induced character at level $N$ the source value is
$\operatorname{den}(N,k)B_{k,\eta}$; it has no additional Euler
factors at primes dividing $N/f$.

% BEGIN OS7
\begin{proposition}[OS7: finite theta evaluation]
Let $e\ge1$, $k=e+1$, and let $\eta$ be primitive of conductor
$f$, with values in $O$, satisfying
$\eta(-1)=(-1)^k$.  Put
\[
       \alpha=\chi_p^{-e}\eta^{-1},\qquad
       Y=B_{k,\eta}/k.
 \tag{OS.33}
\]
Suppose $Y\ne0$ and $\alpha(g_0)-\chi_p(g_0)$ is a unit for
some $g_0$.  For every $a_0\in O\setminus\{0\}$, there are
$m\ge1$, $F=f/p^{v_p(f)}\cdot p^m$, $1\le b<F$ and
representatives $\sigma_t$ such that (OS.9) holds and the element
$z$ in (OS.16) satisfies
\[
              v(z)=v(Y),\qquad a_0z\mid p^m.
 \tag{OS.34}
\]
In particular, if $f\mid N$ and
$a=\operatorname{den}(N,k)B_{k,\eta}\in O$, then $z\mid a$.
The field, the coefficient ring and the character are unchanged.
\end{proposition}
\begin{proof}
We use a simple denominator-clearing version of the source's
Bernoulli congruence.  Choose a nonzero integer $D=D_k$ clearing
the finitely many rational coefficients that occur below.
It depends only on $k$, and not on $F,b,t$ or the character.
Writing
\[
 b_{F,k}(t)=F^{k-1}B_k(t/F)
          =\frac{t^k}{F}-\frac{k}{2}t^{k-1}
             +\sum_{j=2}^k\binom{k}{j}B_jF^{j-1}t^{k-j},
\]
take $u=\langle t\rangle_F$, $y=\langle bt\rangle_F$,
and $h=\lfloor bu/F\rfloor$, with $b$ a unit modulo $F$.
Then $y=bu-Fh$ and
$\theta_{F,b}(t)=2h-(b-1)$.
Expand $y^k$ modulo $F^2$ and $y^{k-1}$ modulo $F$.
After the leading terms cancel, this gives
\[
 D\left[b^{k-1}\theta_{F,b}(t)u^{k-1}
       +2\left(\frac{b_{F,k}(y)}k
                   -b^k\frac{b_{F,k}(u)}k\right)\right]\in F\mathbb Z.
 \tag{OS.35}
\]
This calculation also explains why a denominator independent of
the increasing level suffices: all terms after the two displayed
leading Bernoulli terms contain a factor $F$.
The source uses the explicit choice
$D=2k\prod_{\ell\le k+1}\ell$; the argument here only needs a
fixed common denominator and therefore does not invoke a theorem
on Bernoulli denominators.

Let $\eta_F$ be $\eta$ induced to $F$ and put
\[
 S=\sum_{t\in(\mathbb Z/F)^\times}
          \theta_{F,b}(t)\langle t\rangle_F^{e}\eta_F(t).
\]
Multiply (OS.35) by $\eta_F(bt)$ and sum.
Reindexing the terms containing $y=bt$ gives
\[
 D\left[b^e\eta_F(b)S
       -2\bigl(\eta_F(b)b^k-1\bigr)B_{k,\eta_F}/k\right]
                                               \in p^m O.
 \tag{OS.36}
\]
The generalized Bernoulli numbers here are algebraic numbers in
$K$.  Adding only powers of $p$ to the conductor gives
\[
 B_{k,\eta_F}=
 \begin{cases}
 B_{k,\eta},&p\mid f,\\
 (1-\eta(p)p^{k-1})B_{k,\eta},&p\nmid f,
 \end{cases}
 \tag{OS.37}
\]
once $m\ge v_p(f)$ and $m\ge1$.
For a direct verification, use
$\sum_{i=0}^{q-1}B_k((x+i)/q)=q^{1-k}B_k(x)$,
which follows by summing the generating function
$te^{xt}/(e^t-1)$.  Summation over all lifts at level $qf$
preserves the Bernoulli value; if $q\nmid f$, removing the terms
divisible by $q$ subtracts $\eta(q)q^{k-1}B_{k,\eta}$.
Once $q$ divides the level there is no further removal.
This proves (OS.37).  Its displayed factor is a $p$-adic unit
because $k\ge2$.

Choose $m$ sufficiently large that $f\mid F$ and
\[
                  |p^m|_p<|a_0|_p\,|D|_p\,|Y|_p.
 \tag{OS.38}
\]
Choose $b$ to be the least positive representative of the action of
$g_0$ on $\mu_F$.  Then $b$ is a unit modulo $F$ and
$\eta_F(b)b^k-1$ is a unit in $O$:
modulo the maximal ideal it is associated to
$\alpha(g_0)-\chi_p(g_0)$, after multiplication by the unit
$-\chi_p(g_0)^e\eta(g_0)$.
The factors $b^e\eta_F(b)$ and $2$ are also units.
By (OS.36)--(OS.38) and the ultrametric inequality,
\[
                             |S|_p=|Y|_p.
\]
Indeed the error after division by $D$ is strictly smaller than
the norm of the nonzero term
$2(\eta_F(b)b^k-1)B_{k,\eta_F}/k$.

For any lift $\sigma_t$ of $t$ we have
$\alpha(\sigma_t)^{-1}=\chi_p(\sigma_t)^e\eta_F(t)$ and
$\chi_p(\sigma_t)\equiv\langle t\rangle_F\pmod{p^m}$.
Therefore $z-S\in p^mO$.
By (OS.38) again, $|z|_p=|S|_p=|Y|_p$.
In particular $Y$ has nonnegative valuation, as $z\in O$.
The same inequality implies $v(p^m)\ge v(a_0)+v(z)$, proving
the divisibility in (OS.34).
Also
\[
        a=\bigl(k\operatorname{den}(N,k)\bigr)Y
\]
has valuation at least $v(Y)=v(z)$, so $z\mid a$.
This last argument avoids assuming separately that the primitive
Bernoulli normalization is integral in $O$.

Finally $\alpha|_U\equiv1\pmod{p^m}$ since $f\mid F$,
and parity gives
$\alpha(j)=(-1)^{-e}(-1)^{-(e+1)}=-1$.
The remaining condition in (OS.9) was the hypothesis on $g_0$.
All constructions took place in the original $\Omega$ and $O$.
\end{proof}
% END OS7

\subsection{Normalized Selmer classes and the residual exception}

% BEGIN OS8
\begin{proposition}[OS8: the full cohomology module is controlled]
Let $\eta$ be a primitive Dirichlet character of conductor $f$,
of order prime to $p$, and put $\alpha=\chi_p^{-e}\eta^{-1}$,
with the parity and $e\ge1$ of OS7.
Every class of $\mathscr S(\alpha)$ has a unique cocycle representative
normalized by $c(j)=0$, and these representatives are the increasing
union of $\mathcal T(fp,\varpi^h)$, $h\ge0$.
Consequently a uniform bound $L$ for their lengths gives
$\operatorname{length}_O\mathscr S(\alpha)\le L$.
If
\[
                   \alpha(g)\equiv\chi_p(g)\pmod\varpi
                      \quad\hbox{for every }g\in G,
 \tag{OS.39}
\]
then $\mathscr S(\alpha)=0$.
\end{proposition}
\begin{proof}
Since $\alpha(j)-1=-2$ is a unit, subtracting the coboundary of
$(\alpha(j)-1)^{-1}c(j)$ gives a representative with $c(j)=0$.
Two such representatives differ by a coboundary
$(\alpha(g)-1)m$ vanishing at $j$; this forces $m=0$.
Normalization is therefore unique and $O$-linear.

We check carefully the local conditions in the asserted union.
The character $\alpha$ is unramified outside $fp$.
At each prime $q\mid f$, $q\ne p$, there is
$\tau\in I_q$ with $\alpha(\tau)-1$ a unit.
In fact primitivity means $\eta$ is nontrivial on the $q$-primary
factor of $(\mathbb Z/f)^\times$; otherwise it would factor
through a proper divisor of $f$.  Cyclotomic inertia surjects onto
that factor.  A nontrivial root of unity of order prime to $p$
cannot be congruent to one modulo $\varpi$: if $u^n=1$,
$(n,p)=1$ and $u\ne1$, the identity
$1+u+\cdots+u^{n-1}=0$ would reduce to $n=0$ in $k$.
As $\chi_p$ is unramified at $q\ne p$, this proves the unit claim.

Take a normalized cocycle representing a class in (OS.1).
Its finite image is killed by some $\varpi^h$, since a continuous
map from the compact group $G$ to the discrete torsion group $A$
has finite image.  Choose also one local coboundary witness at
$p$, and enlarge $h$ to kill it.
Conjugating that witness to every other place above $p$ replaces
it by $\alpha(g)m-c(g)$, so the same $\varpi^h$ kills all those
witnesses.
Away from $fp$ the inertia character is trivial; local cohomological
triviality therefore says that the actual cocycle is zero on inertia.
At $q\mid f$, $q\ne p$, let $m_q$ be an inertia coboundary
witness.  The unit value $\alpha(\tau)-1$ just found shows that
$\varpi^h m_q=0$.  It also promotes the witness from inertia to
the decomposition group: for $d\in D_q$, compare the cocycle
identities for $\tau$ and $d^{-1}\tau d$.
They give
\[
  (\alpha(\tau)-1)
       \bigl(c(d)-(\alpha(d)-1)m_q\bigr)=0,
\]
and the coefficient is a unit.  Thus the normalized cocycle belongs
to $\mathcal T(fp,\varpi^h)$.
Conversely the defining conditions of this module imply all the
local triviality conditions in (OS.1), proving the union assertion.

Module length respects this increasing union without a noetherian
assumption.  For any finite strictly increasing chain of submodules
of the union, choose one element witnessing each strict inclusion.
All the finitely many elements occur in one stage; intersecting the
chain with that stage leaves all its inclusions strict.
A common finite length bound for the stages therefore bounds every
such chain and hence the length of the union.

It remains to treat (OS.39), which the finite derivative argument
deliberately excluded.  At a prime $q\ne p$, cyclotomic inertia is
trivial.  If $\eta$ were nontrivial on that inertia group, the
root-of-unity unit argument above would contradict (OS.39).
Thus $\alpha$ is unramified outside $p$.
Moreover $A(\alpha)^G=0$, because $\alpha(j)-1=-2$ is a unit.
Also $A(\alpha)^{D_p}=0$: inertia at $p$ surjects onto
$\operatorname{Gal}(\mathbb Q(\zeta_p)/\mathbb Q)$, so it contains
an element with $\chi_p\equiv-1\pmod\varpi$; by (OS.39) its
$\alpha-1$ is a unit.

Apply cohomology to
$0\to A[\varpi](\alpha)\to A(\alpha)
\xrightarrow{\varpi}A(\alpha)\to0$.
The two vanishing invariant groups just established imply that a
$\varpi$-torsion class in $\mathscr S(\alpha)$ lifts uniquely to
$H^1(\mathbb Q,A[\varpi](\alpha))$ and is still locally zero at
$p$.  At inertia away from $p$, the action on $A$ is trivial and
multiplication by $\varpi$ on its invariants $A$ is surjective;
the same exact sequence makes the lifted class unramified there.
By (OS.39), $A[\varpi](\alpha)\simeq k(1)$.
Choose an $\mathbb F_p$-basis of $k$ and apply the Kummer
identification to each component.  Each component is represented
by $t\in\mathbb Q^\times$.
Unramifiedness at $q\ne p$ forces $v_q(t)\equiv0\pmod p$:
a $p$th root of a unit is unramified there, while a valuation not
divisible by $p$ forces ramification of order $p$.
Local triviality at $p$ makes $t$ a $p$th power in $\mathbb Q_p$,
so $v_p(t)$ is divisible by $p$ as well.
All rational prime exponents in $t$ are thus divisible by $p$,
and its sign is a $p$th power since $p$ is odd.
Hence $t\in\mathbb Q^{\times p}$ and every lifted component is zero.
We have proved $\mathscr S(\alpha)[\varpi]=0$.
Every cohomology class here is killed by a power of $\varpi$,
again because its cocycle has finite image.  Any nonzero such class
would yield nonzero $\varpi$-torsion by multiplying by its largest
nonvanishing power.  Therefore $\mathscr S(\alpha)=0$, as claimed.
\end{proof}
% END OS8

\subsection{Closing the original-field and eigenspace interface}

% BEGIN OS9
\begin{theorem}[OS9: the literal \texttt{PsiBound} supplier]
Let $E/\mathbb Q$ be the original abelian CM field, with the
specified scalar towers into $\Omega$, and let
$\rho:(\mathbb Z/N)^\times\to\operatorname{Gal}(E/\mathbb Q)$,
$N>0$, have the specified cyclotomic compatibility
\texttt{IsCycAct}.  Let $p$ be odd with $p\nmid[E:\mathbb Q]$,
let $O\subset\overline{\mathbb Q}_p$ be the specified finite DVR,
and let $e\ge1$ and
$\psi:\operatorname{Gal}(E/\mathbb Q)\to O^\times$ satisfy
\[
                 (\psi\circ\rho)(-1)=(-1)^{e+1}.
\]
For every supplied $a\in O$ with
\[
           a=\operatorname{den}(N,e+1)
                 B_{e+1,(\psi\circ\rho)_{\rm prim}},
 \tag{OS.40}
\]
the original eigenspace length satisfies
\[
                         \operatorname{psiLen}(\Omega,e,\psi)
                                  \le v_O(a).
 \tag{OS.41}
\]
This is the universal source proposition \texttt{PsiBound}, with
all of its original fields, coefficient rings and character readings.
\end{theorem}
\begin{proof}
Let $\eta_N$ be the Dirichlet character obtained from
$\psi\circ\rho$, and let $\eta$ be its primitive inducer of
conductor $f\mid N$.  Its values lie in the same $O$:
reduction of unit groups onto $(\mathbb Z/f)^\times$ is surjective,
and the induced values are independent of a lift.
Its order divides $|\operatorname{Gal}(E/\mathbb Q)|=[E:\mathbb Q]$,
so is prime to $p$.  Parity is unchanged.
The Galois character attached to $\eta_N$ is identical to that
attached to $\eta$, since their cyclotomic reductions commute.
Write this character as $\widetilde\eta$ and put
$\alpha=\chi_p^{-e}\widetilde\eta^{-1}$.

If $a=0$, (OS.41) is immediate because its right side is $\infty$.
Otherwise the Bernoulli value $Y=B_{e+1,\eta}/(e+1)$ is nonzero.
If (OS.39) holds, OS8 makes $\mathscr S(\alpha)=0$.
If it does not hold, there is a $g_0$ with
$\alpha(g_0)-\chi_p(g_0)$ a unit, since $O$ is local.
For each $a_0\ne0$, apply OS7, then OS6, to obtain
\[
 \operatorname{length}\mathcal T(F,a_0)
                    \le v(z)\le v(a).
\]
All these levels $F=f/p^{v_p(f)}\cdot p^m$ have exactly the
same prime divisors as $fp$.  The definition of $\mathcal T$
depends on this prime support, so
$\mathcal T(F,a_0)=\mathcal T(fp,a_0)$.
OS8 now gives, in both cases,
\[
                   \operatorname{length}_O\mathscr S(\alpha)
                                      \le v_O(a).
 \tag{OS.42}
\]
This also explains a useful simplification: Bernoulli nonvanishing
need not be imported into this supplier.  Its zero case was already
covered by the extended-valuation convention.  Separate nonvanishing
is used later when a finite quantitative bound is wanted.

We finally identify (OS.42) with the exact group in (OS.41).
Cyclotomic compatibility of $\rho$ says, for every $g\in G$,
\[
                  \rho(g|_{\mu_N})=g|_E.
 \tag{OS.43}
\]
Indeed both sides have the same action on every element of the
given copy of $E$ in $\Omega$, by \texttt{IsCycAct} and the
definition of restriction.  Hence
$\widetilde\eta=\psi\circ\operatorname{res}_E$ as actual
characters on $G$, not merely as isomorphic representations.
Let $U_E=\operatorname{Gal}(\Omega/E)$ and
$\Delta=G/U_E=\operatorname{Gal}(E/\mathbb Q)$.
Its order is a unit in $O$.

Restriction in cohomology therefore gives
\[
 H^1(G,A(\alpha))\simeq H^1(U_E,A(\chi_p^{-e}))^{\psi}.
 \tag{OS.44}
\]
Here is the elementary descent argument including the twist sign.
For a finite group of unit order, positive-degree cohomology
vanishes by the averaging homotopy; inflation--restriction is
therefore an isomorphism onto invariants.
On $U_E$ the finite character is trivial, so the coefficient
module really is the same $A(\chi_p^{-e})$.
The ordinary action on a restricted cocycle is
$z(u)\mapsto\chi_p(g)^{-e}z(g^{-1}ug)$.
The invariant action for the twisted coefficient $\alpha$ instead
uses $\alpha(g)=\chi_p(g)^{-e}\psi(g|_E)^{-1}$.
Multiplying by $\psi(g|_E)$ turns this invariant condition into
the $\psi$-eigencondition in (OS.44).

The isomorphism preserves the exact local kernels.
For each decomposition group and each inertia group of $G$,
its intersection with $U_E$ has quotient a subgroup of $\Delta$,
again of unit order.  The same inflation--restriction argument
makes local restriction injective.  Thus a global class is locally
trivial before restriction if and only if it is trivial on all the
corresponding local groups after restriction.
Primes above $p$ correspond to the specified decomposition groups;
all finite inertias correspond to the other part of (OS.1).
Conjugation simply permutes these actual places, so this argument
works for every original choice of them.
Consequently (OS.44) restricts to an $O$-linear isomorphism
\[
 \mathscr S(\alpha)\simeq
 \operatorname{Sha}(E,A(-e))^{\psi},
 \tag{OS.45}
\]
whose right side is precisely the module defining
$\operatorname{psiLen}(\Omega,e,\psi)$.
Length is invariant under this isomorphism.  Combining (OS.42)
and (OS.45) proves (OS.41), with every required quantifier retained.
\end{proof}
% END OS9

\paragraph{Closed dependency boundary.}
OS1--OS9 supply \texttt{PsiBound}; they do not assume it at an
intermediate stage.  Their only duality boundary is LD2--LD4,
whose classical inputs are stated there.  Their remaining classical
arithmetic inputs are the finite cyclotomic/local facts, Hilbert
class field, Chebotarev and Kummer theorem stated at the start.
The Jacobi identities, roots-of-unity removal, derivative construction,
local transition, simultaneous prime selection, finite-length descent,
exceptional residual case and tame eigenspace transport are proved
here.  For the classical number-theoretic inputs see J.~S.~Milne,
\href{https://www.jmilne.org/math/CourseNotes/ANT.pdf}{\emph{Algebraic
Number Theory}, version 3.08}, Sections~6--7 and Theorem~8.31,
and
\href{https://www.jmilne.org/math/CourseNotes/CFT.pdf}{\emph{Class
Field Theory}, version 4.03}, V, Theorem~3.6, Corollary~3.7 and
Example~3.9, and II, Proposition~1.22 for Hilbert's theorem~90.
These references specify general classical theorem leaves; the
component provenance records the exact source declarations and pinned
file hashes for the reconstructed special-value argument.

The literal output (OS.41) is the input to the following odd-minus
quantitative argument.  Thus the former bespoke \texttt{PsiBound}
frontier can be removed once the stated classical inputs and LD are
adopted; this does not remove the separate odd-plus and two-primary
cohomological suppliers or the geometric common-family obligations.
